\section{Follow a checksum from bytes to a file}
The previous project changed how readings were represented. This project leaves the
bytes unchanged and computes a small value from them. Save that value with the data,
compute it again after reading, and compare the two results. A mismatch tells the
reader that the checked bytes and saved value do not agree. Agreement is useful for
finding accidental damage; it does not prove that a message came from a trusted sender.

We will first process an array already in simulator memory. That makes every byte and
state transition visible without file handling. Next, we will replace eight repeated
bit steps with a lookup table. Finally, a file reader will call the same update routine
on each block returned by the operating environment. These are three implementations
of different responsibilities: the checksum state, an acceleration of its update, and
the movement of bytes into a bounded buffer.

By the end, the program should calculate the same checksum for a file read one byte
at a time or in larger blocks. That equality is an interface requirement. The location
of a read boundary must not change the bytes being checked or restart the checksum.
A positive short read is therefore ordinary input to the algorithm, not evidence
that the file has ended.
